\subsection{豪斯多夫群}

设 \(\mathbf{X}\) 是一个集合。使用 § 5 第 1、2 号的记号。自由李代数 \(\mathrm{L}(\mathbf{X})\) 被认作其在 \(\mathrm{A}(\mathbf{X})\) 中的典范像（§ 3，第 1 号，定理 1）。记 \(\hat{\mathrm{L}}(\mathbf{X})\) 为 \(\mathrm{L}(\mathbf{X})\) 在 \(\hat{\mathrm{A}}(\mathbf{X})\) 中的闭包，即 \(\hat{\mathrm{A}}(\mathbf{X})\) 中形如 \(a = \sum_{n\geqslant 1}a_n\) 且满足对所有 \(a_{n}\in \mathrm{L}^{n}(\mathrm{X})\) 都有 \(n\geqslant 0\) 的元素所成的集合；这是 \(\hat{\mathrm{A}}(\mathbf{X})\) 的滤李子代数。

定理 1。将 \(\hat{\mathbf{A}}(\mathbf{X})\) 的指数映射限制在 \(\hat{\mathbf{L}}(\mathbf{X})\) 上，得到从 \(\hat{\mathbf{L}}(\mathbf{X})\) 到 Magnus 群 \(\Gamma(\mathbf{X})\) 的一个闭子群的双射。

记 \(\mathrm{A}(\mathrm{X}) = \mathrm{A}\)、\(\mathrm{A}^{\mathrm{n}}(\mathrm{X}) = \mathrm{A}^{\mathrm{n}}\)、\(\hat{\mathrm{A}}(\mathrm{X}) = \hat{\mathrm{A}}\)、\(\mathrm{L}^{\mathrm{n}}(\mathrm{X}) = \mathrm{L}^{\mathrm{n}}\)、\(\hat{\mathrm{L}}(\mathrm{X}) = \hat{\mathrm{L}}\)、\(\Gamma(\mathrm{X}) = \Gamma\)。令 \(\mathrm{B}\) 为代数 \(\mathrm{A} \otimes \mathrm{A}\)，其带有由 \(\mathbf{N}\) 定义的 \(\mathrm{B}^{\mathrm{n}} = \sum_{i+j=n} \mathrm{A}^{i} \otimes \mathrm{A}^{j}\) 型分次。令 \(\hat{\mathrm{B}} = \prod_{n > 0} \mathrm{B}^{n}\) 为相应的完备滤代数（《交换代数》，第 III 章，§2，第 12 号，例 1）。§3 第 1 号定理 1 的推论 1 中定义的余乘法 \(c: \mathrm{A} \to \mathrm{A} \otimes \mathrm{A}\) 是 0 次齐次的，因而连续延拓为同态 \(\hat c: \hat{\mathrm{A}} \to \hat{\mathrm{B}}\)，其给出为

\[
\hat {c} \left(\sum_ {n \geqslant 0} a _ {n}\right) = \sum_ {n \geqslant 0} c (a _ {n}) \quad \text { for } a _ {n} \in \mathrm{A} ^ {n}.
\]

我们还定义从 \(\delta'\) 到 \(\delta''\) 的连续同态 \(\hat{\mathbf{A}}\) 和 \(\hat{\mathbf{B}}\)，规定为

\[
\delta^ {\prime} \left(\sum_ {n \geqslant 0} a _ {n}\right) = \sum_ {n \geqslant 0} (a _ {n} \otimes 1), \quad \delta^ {\prime \prime} \left(\sum_ {n \geqslant 0} a _ {n}\right) = \sum_ {n \geqslant 0} (1 \otimes a _ {n}) \quad \text { for } a _ {n} \in \mathrm{A} ^ {n}.
\]

根据 §3 第 1 号定理 1 的推论 2，\(\mathbf{L}^n\) 是满足 \(a_{n}\in \mathbf{A}^{n}\) 的 \(c(a_{n}) = a_{n}\otimes 1 + 1\otimes a_{n}\) 所成的集合。由此，\(\hat{\mathbf{L}}\) 是满足下式的 \(a\in \hat{\mathbf{A}}\) 所成的集合：

\[
\hat {c} (a) = \delta^ {\prime} (a) + \delta^ {\prime \prime} (a).\tag{3}
\]

令 \(\Delta\) 为 \(b\in \hat{\mathbf{A}}\) 中常数项等于 1 且满足关系

\[
\hat {c} (b) = \delta^ {\prime} (b). \delta^ {\prime \prime} (b),\tag{4}
\]

换言之，是满足 \(b = \sum_{n\geqslant 0}b_n\)、\(b_{n}\in \mathbf{A}^{n}\)（对所有 \(n\geqslant 0\)）、\(b_0 = 1\) 且满足 \(c(b_{n}) = \sum_{t + j = n}b_{t}\otimes b_{j}\)（对 \(n\geqslant 0\)）的 b 所成的集合。后一个刻画表明 \(\Delta\) 是 \(\Gamma\) 的闭子集；由于 \(\hat c,\delta^{\prime}\) 和 \(\delta ''\) 是环同态，且 \(\delta'(\hat{\mathbf A})\) 的每个元素都与 \(\delta''(\hat{\mathbf A})\) 的每个元素可交换，在 \(\Gamma\) 上，映射 \(c\) 和 \(\delta ' \delta''\) 的限制是群同态，而 \(\Delta\) 是 \(\Gamma\) 的子群。

由第 1 号命题 1，\(\hat{\mathbf{A}}\) 的指数映射是从无常数项的 \(\hat{\mathbf{A}}^{+}\) 元素集 \(\hat{\mathbf{A}}\) 到 \(\Gamma\) 的双射。令 \(a \in \hat{\mathbf{A}}^{+}\) 且 \(b = \exp a\)。由于 \(\hat c\) 是连续环同态，

\[
\hat {c} (b) = \hat {c} \left(\sum_ {n \geqslant 0} a ^ {n} / n!\right) = \sum_ {n \geqslant 0} \hat {c} (a) ^ {n} / n! = \exp \hat {c} (a).
\]

关系

\[
\delta^ {\prime} (b) = \exp \delta^ {\prime} (a), \quad \delta^ {\prime \prime} (b) = \exp \delta^ {\prime \prime} (a)
\]

类似可证；又由于 \(\delta'(a)\) 与 \(\delta''(a)\) 可交换（第 1 号，注 1），

\[
\delta^ {\prime} (b) \delta^ {\prime \prime} (b) = \exp (\delta^ {\prime} (a) + \delta^ {\prime \prime} (a)).
\]

因此，a 满足 (3) 当且仅当 b 满足 (4)，定理得证。注。上述证明表明，\(\exp(\hat{\mathbf{L}})\) 是 \(\Delta\) 的子群 \(\Gamma\)，它由满足 (4) 的 b 构成。

因此，\(\Delta\) 的群律可通过指数映射传递到 \(\hat{\mathbf{L}}\)。换言之，\(\hat{\mathbf{L}}\) 是一个完备拓扑群，其复合律 \((a, b) \mapsto a \uplus b\) 由下式给出

\[
a \uplus b = \log (\exp a. \exp b).
\]

如此得到的拓扑群称为豪斯多夫群（由 X 关于 K 导出）。

令 \(g\) 为从自由群 \(\mathbf{F} = \mathrm{F}(\mathbf{X})\) 到 \(\Gamma\) 的同态，满足对 \(g(x) = \exp x\) 有 \(x \in \mathbf{X}\)。由于 \(\exp x - 1 - x = \sum_{n \geqslant 2} x^n / n!\) 的阶 \(\geqslant 2\)，由 §5 第 3 号定理 1，\(g\) 是单射。因此，映射 \(\log \circ g\) 是 \(F\) 到豪斯多夫群的单射同态，并延拓了典范嵌入 \(X \to L\)。

对每个整数 \(m \geqslant 1\)，记 \(\hat{\mathbf{L}}_m\) 为 \(\geqslant m\) 中阶 \(\hat{\mathbf{L}}\) 的元素所成的集合，记 \(\Gamma_m\) 为满足 \(u \in \Gamma\) 的阶 \(u - 1\) 的 \(\geqslant m\) 所成的集合。由第 1 号的注 2，\(\hat{\mathbf{L}}_m = \exp^{-1}(\Gamma_m)\)；由于 \((\Gamma_m)_{m \geqslant 1}\) 是 \(\Gamma\) 上的整中心滤子（§4，第 5 号，命题 2），\((\hat{\mathbf{L}}_m)_{m \geqslant 1}\) 是群 \(\hat{\mathbf{L}}\) 上的整中心滤子。

